\documentclass[12pt]{article}
\usepackage[utf8]{inputenc}
\usepackage[T1]{fontenc}
\usepackage{hyperref}
\title{Wogeo Language Resources for K-12 Curriculum Development}
\author{Compiled by Claude Sonnet 5.5 (claude-sonnet-5-5) for Andr\'as Kornai}
\date{}
\begin{document}
\maketitle

\section{Introduction}
Wogeo (ISO 639-3 woc, Glottocode woge1237) is an Oceanic (Austronesian) language of the Schouten branch, in the Kairiru--Manam group. It is spoken on Koil and Vokeo islands in East Sepik Province, Papua New Guinea. Wikipedia reports about 1{,}600 speakers (2003). Search found no school materials.

\section{Primary Educational Resources}
No primers, textbooks or school programmes were found. The language is not known to be taught in schools.

\section{Linguistic and Reference Materials}
Wikipedia has a language article with a phoneme inventory: \url{https://en.wikipedia.org/wiki/Wogeo_language}.

A Wogeo document appears on the DICE (Missouri) Austronesian site: \url{https://dice.missouri.edu/assets/docs/austronesia/Wogeo.pdf}. I did not read it, so its content is unconfirmed.

Encyclopedia.com has an ethnographic entry on the Wogeo people, who call themselves Wageva: \url{https://www.encyclopedia.com/humanities/encyclopedias-almanacs-transcripts-and-maps/wogeo}.

\section{Religious and Community Resources}
Joshua Project reports no Scripture in Wogeo, and mother-tongue literacy below 5\%: \url{https://joshuaproject.net/people_groups/15406}. It suggests that a translation might draw on the related Manam New Testament.

\section{Digital and Audio Resources}
None found.

\section{Curriculum Development Context}
Wogeo is still spoken across generations, but with decreasing use among young people. Literacy is low and there is no known orthography-based material. A curriculum would start from basic documentation and a primer, with the community.
\end{document}